Este capítulo constituye la parte principal del libro. De cierto modo, el material de los otros capítulos es preparatorio, auxiliar o adicional en relación a métodos para la resolución de problemas de optimización con restricciones.

\section{Métodos para conjuntos factibles de estructura simple}

Comenzaremos el desarrollo de métodos para problemas de optimización con restricciones en el caso en que el conjunto factible tenga descripción abstracta:
\begin{equation}\label{eq:4.1}
    \min f(x) \quad \text{sujeto a} \quad x \in D,
\end{equation}
donde \( \inmap{f}{\R^n}{\R} \) es una función diferenciable y $D \subset \R^n$. Por razones de aplicación práctica, entendemos que el conjunto $D$ posee una estructura simple (por lo menos, $D$ es convexo y cerrado).

Con frecuencia, problemas con restricciones simples aparecen como subproblemas en métodos para la resolución de problemas con restricciones más complejas. Notamos que cuando el conjunto $D$ es convexo y cerrado, para todo punto $x \in \R^n$ existe una proyección $P_D(x)$ de $x$ sobre $D$ que es única (vea el Teorema 1.3.1). Recordamos también que el punto $\bar{x} \in D$ es un punto estacionario del problema \eqref{eq:4.1} cuando
\begin{equation}\label{eq:4.2}
    \interno{f'(\bar{x})}{x - \bar{x}} \ge 0 \quad \forall x \in D,
\end{equation}
o, equivalentemente,
\begin{equation}\label{eq:4.3}
    P_D(\bar{x} - \alpha f'(\bar{x})) = \bar{x}
\end{equation}
para algún $\alpha > 0$ (vea el Teorema 1.3.7). Más aún, la validez de \eqref{eq:4.3} para algún $\alpha > 0$ implica la validez de esta relación para todo $\alpha > 0$.

\subsection{Métodos del gradiente proyectado}

Los métodos del gradiente proyectado pueden ser pensados como una combinación natural de métodos del gradiente para optimización irrestricta con proyección de iteraciones así obtenidas en el conjunto factible del problema. Esto resulta en el esquema siguiente:
\begin{equation}\label{eq:4.4}
    x^{k+1} = P_D(x^k - \alpha_k f'(x^k)), \quad k = 0, 1, \dots,
\end{equation}
donde $x^0 \in D$, y los parámetros de longitud de paso $\alpha_k > 0$ son calculados utilizando extensiones de técnicas de búsqueda lineal de § 3.1.1 para el caso con restricciones. En particular, el valor de longitud de paso es calculado para obtener decrecimiento de la función
\[
    \inmap{\varphi_k}{\R_+}{\R}, \quad \varphi_k(\alpha) = f(x^k(\alpha)),
\]
donde
\begin{equation}\label{eq:4.5}
    x^k(\alpha) = P_D(x^k - \alpha f'(x^k))
\end{equation}
es el arco de proyección (vea Figura 4.1.1).

Por ejemplo, en el caso de la minimización unidimensional, $\alpha_k$ es una solución del problema
\[
    \min \varphi_k(\alpha) \quad \text{sujeto a} \quad \alpha \in \R_+,
\]
o del problema
\[
    \min \varphi_k(\alpha) \quad \text{sujeto a} \quad \alpha \in [0, \hat{\alpha}],
\]
donde $\hat{\alpha} > 0$ es un parámetro. Sin embargo, como en el caso irrestricto, las reglas de minimización unidimensional no son muy atractivas desde el punto de vista práctico.

Más interesantes son las extensiones de la regla de Armijo. Esto puede ser hecho, por ejemplo, por el mismo esquema que fue considerado en § 3.1.1, sustituyendo la condición (3.6) por
\begin{equation}\label{eq:4.6}
    f(x^k(\alpha)) \le f(x^k) + \sigma \interno{f'(x^k)}{x^k(\alpha) - x^k}.
\end{equation}
Recordamos que, para comenzar el proceso, deben ser elegidos un parámetro $\sigma \in (0, 1)$, el valor inicial de la longitud de paso $\hat{\alpha} > 0$ y el parámetro de reducción del paso $\theta \in (0, 1)$. En particular, $\alpha_k$ es calculado como el mayor entre todos los números de la forma $\alpha = \hat{\alpha} \theta^s$, $s \in \N$, satisfaciendo la desigualdad \eqref{eq:4.6}. Como es fácil ver, en el caso cuando $D = \R^n$, la desigualdad \eqref{eq:4.6} se reduce a la condición de Armijo (3.6) en el caso irrestricto.

Existen también varias modificaciones de la regla de Armijo (por ejemplo, similares a la regla de Goldstein). Podríamos considerar también la regla de longitud de paso fijo, donde $\alpha_k = \bar{\alpha}$ para todo $k$, con $\bar{\alpha} > 0$ fijo.

\begin{algorithm}[Método del gradiente proyectado]
    Elegir $x^0 \in D$ y tomar $k := 0$. Elegir una de las tres reglas de cálculo de longitud de paso a ser utilizada, y los parámetros asociados: $\hat{\alpha} > 0$ (o $\hat{\alpha} = +\infty$) en el caso de minimización unidimensional, $\hat{\alpha} > 0$, $\sigma, \theta \in (0, 1)$ en el caso de la regla de Armijo, $\bar{\alpha} > 0$ en el caso de longitud de paso fijo.
    
    \begin{enumerate}
        \item Calcular $\alpha_k$ de acuerdo con la regla elegida.
        \item Calcular $x^{k+1}$ utilizando la fórmula \eqref{eq:4.4}.
        \item Tomar $k := k + 1$ y retornar al Paso 1.
    \end{enumerate}
\end{algorithm}

El Algoritmo 4.1.1 está ilustrado en la Figura 4.1.1.

Observamos que si para algún $k$ el punto $x^k$ es estacionario para el problema \eqref{eq:4.1}, entonces el método genera $x^k = x^{k+1} = \dots$, cualquiera que sea la longitud de paso $\alpha_k$. Notamos también que para $D = \R^n$, un método del gradiente proyectado se reduce al método del gradiente (irrestricto) correspondiente. Las propiedades de convergencia de métodos del gradiente proyectado son muy parecidas a las propiedades de métodos del gradiente. A continuación, presentaremos dos resultados que son generalizaciones de los Teoremas 3.1.1 y 3.1.5, respectivamente.

En el análisis a seguir, la siguiente desigualdad tiene un papel importante:
\begin{equation}\label{eq:4.7}
    \interno{f'(x^k)}{x^k(\alpha) - x^k} \le -\frac{1}{\alpha} \| x^k(\alpha) - x^k \|^2 \quad \forall \alpha > 0.
\end{equation}
Para probar esta desigualdad, notamos que para todo $\alpha > 0$,
\begin{align*}
    \interno{f'(x^k)}{x^k(\alpha) - x^k} &= \frac{1}{\alpha} \interno{x^k + \alpha f'(x^k) - x^k}{x^k(\alpha) - x^k} \\
    &= \frac{1}{\alpha} \interno{x^k - x^k(\alpha)}{x^k(\alpha) - x^k} + \frac{1}{\alpha} \interno{x^k(\alpha) - (x^k - \alpha f'(x^k))}{x^k(\alpha) - x^k} \\
    &\le -\frac{1}{\alpha} \| x^k(\alpha) - x^k \|^2,
\end{align*}
donde el Teorema 1.3.1 fue utilizado.

\begin{lemma}
    Sean $D \subset \R^n$ un conjunto cerrado y convexo, y \( \inmap{f}{\R^n}{\R} \)una función diferenciable en $\R^n$, con derivada Lipschitz-continua en $\R^n$ con módulo $L > 0$.
    
    Entonces para cualquier $x^k \in \R^n$ la desigualdad \eqref{eq:4.6} es satisfecha para todo $\alpha \in (0, \bar{\alpha}]$, donde
    \begin{equation}\label{eq:4.8}
        \bar{\alpha} = 2(1 - \sigma)/L > 0.
    \end{equation}
\end{lemma}

\begin{proof}
    Utilizando el Lema 1.5.4 y las desigualdades \eqref{eq:4.7} y \eqref{eq:4.8}, para todo $\alpha \in (0, \bar{\alpha}]$ obtenemos que
    \begin{align*}
        f(x^k(\alpha)) - f(x^k) &\le \interno{f'(x^k)}{x^k(\alpha) - x^k} + \frac{L}{2} \| x^k(\alpha) - x^k \|^2 \\
        &\le \sigma \interno{f'(x^k)}{x^k(\alpha) - x^k} + \left( \frac{L}{2} - \frac{1 - \sigma}{\alpha} \right) \| x^k(\alpha) - x^k \|^2 \\
        &\le \sigma \interno{f'(x^k)}{x^k(\alpha) - x^k},
    \end{align*}
    lo que es justamente el resultado anunciado.
\end{proof}

El resultado del Lema 4.1.1 implica en que, bajo las hipótesis correspondientes, los valores de longitud de paso calculados por la regla de Armijo son limitados inferiormente por un número que no depende de $k$, i.e.,
\begin{equation}\label{eq:4.9}
    \alpha_k \ge \check{\alpha} > 0 \quad \forall k.
\end{equation}
Observamos también que la utilización de longitud de paso fijo, suponiendo que
\begin{equation}\label{eq:4.10}
    \bar{\alpha} < 2/L,
\end{equation}
es equivalente a la regla de Armijo con $\hat{\alpha} = \bar{\alpha}$ y $\sigma > 0$ suficientemente pequeño. Por lo tanto, este caso no requiere análisis separado.

\begin{theorem}[Convergencia del método del gradiente proyectado]\label{th:4.1.1}
Sean $D \subset \R^n$ un conjunto cerrado y convexo, $f : \R^n \to \R$ una función diferenciable en $\R^n$, con derivada Lipschitz-continua en $\R^n$ con módulo $L > 0$.

Entonces cada punto de acumulación de cualquier secuencia $\{x^k\}$ generada por el Algoritmo 4.1.1 es un punto estacionario del problema \eqref{eq:4.1}. Si un punto de acumulación existe, o si la función $f$ es limitada inferiormente en $D$, entonces para cualquier secuencia limitada $\{t_k\} \subset \R_+$ se tiene que
\begin{equation}\label{eq:4.11}
    \| x^k - P_D(x^k - t_k f'(x^k)) \| \to 0 \quad (k \to \infty).
\end{equation}
\end{theorem}

\begin{proof}
Consideremos el caso de la regla de Armijo. Por \eqref{eq:4.7}, para todo $k$ tenemos que
\begin{align}
    f(x^k) - f(x^{k+1}) &\ge -\sigma \interno{f'(x^k)}{x^{k+1} - x^k} \nonumber \\
    &\ge \frac{\sigma}{\alpha_k} \| x^{k+1} - x^k \|^2 \nonumber \\
    &\ge \frac{\sigma}{\check{\alpha}} \| x^{k+1} - x^k \|^2. \label{eq:4.12}
\end{align}

Si $x^k$ no es un punto estacionario del problema \eqref{eq:4.1}, se tiene que
\[
    \| x^{k+1} - x^k \| = \| P_D(x^k - \alpha_k f'(x^k)) - x^k \| > 0,
\]
porque $\alpha_k > 0$. Por eso, de \eqref{eq:4.12} obtenemos que la secuencia $\{f(x^k)\}$ es decreciente. Supongamos que la secuencia $\{x^k\}$ posea un punto de acumulación $\bar{x} \in \R^n$. En este caso, por el hecho de que la secuencia $\{f(x^k)\}$ es decreciente, obtenemos que ella es limitada inferiormente (por $f(\bar{x})$), y por lo tanto converge (cuando la función $f$ es limitada inferiormente en $\R^n$, esto vale aun cuando $\{x^k\}$ no tiene puntos de acumulación). Concluimos que
\[
    f(x^k) - f(x^{k+1}) \to 0 \quad (k \to \infty),
\]
y, de \eqref{eq:4.12}, obtenemos que
\begin{equation}\label{eq:4.13}
    \| x^{k+1} - x^k \| \to 0 \quad (k \to \infty).
\end{equation}

Sea la secuencia $\{t_k\} \subset \R_+$ limitada superiormente por un número $\hat{t}$. Por el Teorema 1.3.1, utilizando también \eqref{eq:4.9} y \eqref{eq:4.13}, obtenemos que
\begin{align*}
    \| x^k(t_k) - x^k \| &\le \max\{1, t_k/\alpha_k\} \| x^k(\alpha_k) - x^k \| \\
    &\le \max\{1, \hat{t}/\check{\alpha}\} \| x^{k+1} - x^k \| \to 0 \quad (k \to \infty),
\end{align*}
lo que prueba \eqref{eq:4.11}.

Sea $\{x^{k_j}\}$ una subsecuencia de la secuencia $\{x^k\}$ que converge a $\bar{x}$. Pasando al límite cuando $j \to \infty$ en \eqref{eq:4.11}, donde tomamos $t_k = t \ \forall k$, siendo $t > 0$ un número arbitrario, teniendo en cuenta la continuidad del operador de proyección (vea el Teorema 1.3.1), obtenemos \eqref{eq:4.3}, mostrando que $\bar{x}$ es un punto estacionario.

La consideración del caso de utilización de la minimización unidimensional se reduce a la regla de Armijo de la misma manera que en el Teorema 3.1.1.
\end{proof}

Sea $S_0$ el conjunto de puntos estacionarios del problema \eqref{eq:4.1}. Supongamos que valga la condición de separación de superficies de nivel críticas, que tiene la misma forma que en § 3.1.2 (pero con la nueva definición de $S_0$, claro). Supongamos que vale la estimativa
\begin{equation}\label{eq:4.14}
    \| x - P_D(x - f'(x)) \| \ge \gamma \dist(x, S_0) \quad \forall x \in U,
\end{equation}
donde
\begin{equation}\label{eq:4.15}
    U = \{ x \in \R^n \mid \| x - P_D(x - f'(x)) \| < \delta \},
\end{equation}
con constantes $\delta > 0$ y $\gamma > 0$. En el caso de $D = \R^n$, estas condiciones ya fueron consideradas en § 3.1.2. Sin entrar en detalles, mencionamos que ambas son satisfechas si, por ejemplo, $f$ es una función cuadrática, $D$ es un conjunto poliédrico, y $S_0 \neq \emptyset$ (comparar con el Ejercicio 3.1.11).

\begin{theorem}[Convergencia lineal del método del gradiente proyectado]\label{th:4.1.2}
Además de las hipótesis del Teorema 4.1.1, supongamos que el valor óptimo del problema \eqref{eq:4.1} sea finito y valga la condición \eqref{eq:4.14}, con $U$ definido en \eqref{eq:4.15} y algunos $\delta > 0$ y $\gamma > 0$. Supongamos también que valga la condición de separación de superficies de nivel críticas.

Entonces cualquier secuencia $\{x^k\}$ generada por el Algoritmo 4.1.1 converge a un punto estacionario del problema \eqref{eq:4.1}. La tasa de convergencia con respecto a la función es lineal, y con respecto a las variables la tasa es geométrica.
\end{theorem}

La demostración de este teorema puede ser hecha por la modificación relativamente obvia de la demostración del Teorema 3.1.5. Mencionemos apenas que aquí deben ser también utilizados los resultados de los Teoremas 1.3.1 y 1.3.7.

Cada iteración de un método del gradiente proyectado requiere el cálculo de proyecciones sobre el conjunto $D$ (una vez, o más, dependiendo de la regla para la longitud de paso). Esto indica que el método sólo puede ser práctico cuando el cálculo de proyecciones tiene un costo computacional relativamente bajo. Por ejemplo, cuando $D$ es un conjunto poliédrico, el problema de encontrar una proyección es de programación cuadrática (vea § 6.2), lo que es considerado aceptable. Difícilmente puede tener sentido utilizar técnicas de proyección en el caso más general, i.e., de restricciones no lineales. A veces, el cálculo de proyecciones no requiere resolución numérica de problemas de optimización, porque existen fórmulas explícitas.

Existen modificaciones de métodos del gradiente proyectado que requieren, en cada iteración, una proyección apenas. Por ejemplo, en vez del esquema \eqref{eq:4.4} basado en el arco de proyección \eqref{eq:4.5}, podemos utilizar el intervalo
\[
    x^k(\alpha) = (1 - \alpha)x^k + \alpha P_D(x^k - \beta_k f'(x^k)).
\]
Esto resulta en el esquema
\begin{equation}\label{eq:4.16}
    x^{k+1} = (1 - \alpha_k)x^k + \alpha_k P_D(x^k - \beta_k f'(x^k)), \quad k = 0, 1, \dots,
\end{equation}
donde $x^0 \in D$, los parámetros $\beta_k > 0$ pueden ser elegidos más o menos arbitrariamente (desde que $0 < c_1 \le \beta_k \le c_2$), y el parámetro de longitud de paso $\alpha_k$ es calculado por la regla de Armijo \eqref{eq:4.6} con $\hat{\alpha} = 1$. Las reglas de minimización unidimensional pueden ser adaptadas análogamente.

Observamos que si $x^k \in D$ entonces, por la convexidad de $D$,
\[
    d^k = P_D(x^k - \beta_k f'(x^k)) - x^k
\]
es una dirección factible en el punto $x^k$ con respecto al conjunto $D$ y, en particular, $x^k(\alpha) = x^k + \alpha d^k \in D$ para todo $\alpha \in [0, 1]$. Este intervalo está ilustrado en la Figura 4.1.2.

Además de eso, por \eqref{eq:4.7}, $d^k$ es una dirección de descenso de la función $f$ en el punto $x^k$. Tenemos entonces una dirección de descenso que es factible. A partir de eso, es fácil probar que, si la secuencia $\{\beta_k\}$ es limitada inferiormente por un número positivo, este esquema posee convergencia global en el sentido del Teorema 4.1.1. Sin embargo, probar la tasa de convergencia lineal para este esquema ya no es posible, por lo menos bajo las hipótesis usuales.

La importancia teórica de métodos del gradiente proyectado consiste principalmente en el hecho de que varios otros métodos que generan iteraciones factibles y son métodos de descenso, pueden ser interpretados como versiones inexactas o modificadas de algún método del gradiente proyectado.
%%% Local Variables:
%%% mode: LaTeX
%%% TeX-master: "../../../Apuntes-Programacion-No-Lineal"
%%% End:
